% GENERATED FILE. Edit canonical lessons, then run npm run generate:books.
% canonical-source-hash: fnv1a64:be6518daa261728f
% canonical-lessons: TE-C326-saksi, TE-C326-yuddham, TE-C326-sanghatana, TE-C326-dongilincu, TE-C326-paripovu

\chapter{A Witness to the War}
\label{ch:te-a-witness-to-the-war}
\hlchaptermodality{326}

\begin{chapteropening}
\textbf{By the end of this chapter:} \emph{I can say a witness, a war, an incident, to steal and to run away.}

Complete the last lesson of chapter 326: I can say a witness, a war, an incident, to steal and to run away.
\end{chapteropening}

\section[\texorpdfstring{\te{సాక్షి} (sākṣi)}{సాక్షి (sākṣi)}]{\te{సాక్షి} (sākṣi) — a witness}

\label{lesson:TE-C326-saksi}



\begin{quote}
Before the new one: say the Telugu for a watchman, then the Telugu for a milkman.
\end{quote}

\subsection*{You'll want to know: \te{సాక్షి}}
\textbf{\te{సాక్షి}} — \emph{sākṣi} — \textquotedblleft{}a witness\textquotedblright{}.

\begin{quote}
\te{సాక్షి} \te{పోలీసుతో} \te{మాట్లాడాడు}.
\end{quote}

\begin{quote}
\emph{sākṣi pōlīsutō māṭlāḍāḍu.}
\end{quote}

\begin{quote}
The witness spoke with the police.
\end{quote}

Plural: \te{సాక్షులు}, \emph{sākṣulu}.

A \te{ప్రత్యక్ష} \te{సాక్షి} is an eyewitness, one who saw it happen. A whole case can turn on one: it depends on the witness, \te{సాక్షి} \te{మీద} \te{ఆధారపడి} \te{ఉంటుంది}.

\begin{cousinweb}
\textbf{\te{సాక్షి}} (\emph{sākṣi}) is borrowed from Sanskrit \emph{sākṣin}, \textquotedblleft{}a witness\textquotedblright{} (CDIAL 13321). CDIAL takes the Sanskrit word apart as \emph{sa-} \textquotedblleft{}with\textquotedblright{} and \emph{akṣi} \textquotedblleft{}eye\textquotedblright{}: the witness is the one who was there with his own eyes.
\end{cousinweb}

\begin{grammarlens}[title={What you've built}]
One more word for telling what happened.
\end{grammarlens}

\subsection*{Your turn}
\begin{itemize}
\raggedright
  \item \emph{Say it:} \emph{sākṣi}
  \item \emph{Say it:} \emph{sākṣi}, once more
  \item \emph{Say it:} \emph{sākṣi pōlīsutō māṭlāḍāḍu.} — The witness spoke with the police.
  \item \emph{Recall:} say the Telugu for a watchman, then the Telugu for a milkman, then say \emph{sākṣi} again
\end{itemize}

\begin{tcolorbox}[breakable,colback=teal!4,colframe=teal!35!black,title={Before you move on}]
What does \te{సాక్షి} mean? (\textquotedblleft{}A witness\textquotedblright{}.) Say it once more.
\end{tcolorbox}
\section[\texorpdfstring{\te{యుద్ధం} (yuddhaṁ)}{యుద్ధం (yuddhaṁ)}]{\te{యుద్ధం} (yuddhaṁ) — a war; a battle}

\label{lesson:TE-C326-yuddham}



\begin{quote}
Before the new one: say the Telugu for a milkman, then the Telugu for a witness.
\end{quote}

\subsection*{You'll want to know: \te{యుద్ధం}}
\textbf{\te{యుద్ధం}} — \emph{yuddhaṁ} — \textquotedblleft{}a war; a battle\textquotedblright{}.

\begin{quote}
\te{ఆ} \te{యుద్ధం} \te{చాలా} \te{రోజులు} \te{జరిగింది}.
\end{quote}

\begin{quote}
\emph{ā yuddhaṁ cālā rōjulu jarigindi.}
\end{quote}

\begin{quote}
That war went on for many days.
\end{quote}

With \te{చేయు} it makes a verb: \te{యుద్ధం} \te{చేయు}, to fight a war.

\begin{cousinweb}
\textbf{\te{యుద్ధం}} (\emph{yuddhaṁ}) is borrowed from Sanskrit \emph{yuddha}, \textquotedblleft{}fought\textquotedblright{}, and as a noun \textquotedblleft{}a battle\textquotedblright{}, the past participle of the root \emph{yudh}, \textquotedblleft{}to fight\textquotedblright{} (CDIAL 10499). Like many Sanskrit nouns in -a, it ends in -\te{ం} in Telugu.
\end{cousinweb}

\begin{grammarlens}[title={What you've built}]
One more word for telling what happened.
\end{grammarlens}

\subsection*{Your turn}
\begin{itemize}
\raggedright
  \item \emph{Say it:} \emph{yuddhaṁ}
  \item \emph{Say it:} \emph{yuddhaṁ}, once more
  \item \emph{Say it:} \emph{ā yuddhaṁ cālā rōjulu jarigindi.} — That war went on for many days.
  \item \emph{Recall:} say the Telugu for a milkman, then the Telugu for a witness, then say \emph{yuddhaṁ} again
\end{itemize}

\begin{tcolorbox}[breakable,colback=teal!4,colframe=teal!35!black,title={Before you move on}]
What does \te{యుద్ధం} mean? (\textquotedblleft{}A war; a battle\textquotedblright{}.) Say it once more.
\end{tcolorbox}
\section[\texorpdfstring{\te{సంఘటన} (saṅghaṭana)}{సంఘటన (saṅghaṭana)}]{\te{సంఘటన} (saṅghaṭana) — an incident; an event}

\label{lesson:TE-C326-sanghatana}



\begin{quote}
Before the new one: say the Telugu for a witness, then the Telugu for a war.
\end{quote}

\subsection*{You'll want to know: \te{సంఘటన}}
\textbf{\te{సంఘటన}} — \emph{saṅghaṭana} — \textquotedblleft{}an incident; an event\textquotedblright{}.

\begin{quote}
\te{నిన్న} \te{రాత్రి} \te{ఒక} \te{వింత} \te{సంఘటన} \te{జరిగింది}.
\end{quote}

\begin{quote}
\emph{ninna rātri oka vinta saṅghaṭana jarigindi.}
\end{quote}

\begin{quote}
A strange incident happened last night.
\end{quote}

It goes with \te{జరుగు}, to happen: \te{సంఘటన} \te{జరిగింది}, the incident happened.

\begin{cousinweb}
\textbf{\te{సంఘటన}} (\emph{saṅghaṭana}) is borrowed from Sanskrit \emph{saṁghaṭana}, \textquotedblleft{}union, a coming together\textquotedblright{} (Monier-Williams), built from \emph{sam-} \textquotedblleft{}together\textquotedblright{} and \emph{ghaṭ} \textquotedblleft{}to join\textquotedblright{}. In Telugu the things that come together make an event, an incident.
\end{cousinweb}

\begin{grammarlens}[title={What you've built}]
One more word for telling what happened.
\end{grammarlens}

\subsection*{Your turn}
\begin{itemize}
\raggedright
  \item \emph{Say it:} \emph{saṅghaṭana}
  \item \emph{Say it:} \emph{saṅghaṭana}, once more
  \item \emph{Say it:} \emph{ninna rātri oka vinta saṅghaṭana jarigindi.} — A strange incident happened last night.
  \item \emph{Recall:} say the Telugu for a witness, then the Telugu for a war, then say \emph{saṅghaṭana} again
\end{itemize}

\begin{tcolorbox}[breakable,colback=teal!4,colframe=teal!35!black,title={Before you move on}]
What does \te{సంఘటన} mean? (\textquotedblleft{}An incident; an event\textquotedblright{}.) Say it once more.
\end{tcolorbox}
\section[\texorpdfstring{\te{దొంగిలించు} (doṅgiliñcu)}{దొంగిలించు (doṅgiliñcu)}]{\te{దొంగిలించు} (doṅgiliñcu) — to steal}

\label{lesson:TE-C326-dongilincu}



\begin{quote}
Before the new one: say the Telugu for a war, then the Telugu for an incident.
\end{quote}

\subsection*{You'll want to know: \te{దొంగిలించు}}
\textbf{\te{దొంగిలించు}} — \emph{doṅgiliñcu} — \textquotedblleft{}to steal\textquotedblright{}.

\begin{quote}
\te{దొంగ} \te{నా} \te{సంచి} \te{దొంగిలించాడు}.
\end{quote}

\begin{quote}
\emph{doṅga nā sañci doṅgiliñcāḍu.}
\end{quote}

\begin{quote}
The thief stole my bag.
\end{quote}

Past (I …): \te{దొంగిలించాను}, \emph{doṅgiliñcānu}.

It is the verb of \te{దొంగ}, a thief, which you already know.

\begin{cousinweb}
\textbf{\te{దొంగిలించు}} (\emph{doṅgiliñcu}) is native Dravidian: DEDR 2982 lists Telugu \emph{doṅgilincu} \textquotedblleft{}to steal\textquotedblright{} beside \emph{doṅga} \textquotedblleft{}a thief\textquotedblright{}, with Tamil \emph{toṅkaṉ} \textquotedblleft{}a robber\textquotedblright{}.
\end{cousinweb}

\begin{grammarlens}[title={What you've built}]
One more word for telling what happened.
\end{grammarlens}

\subsection*{Your turn}
\begin{itemize}
\raggedright
  \item \emph{Say it:} \emph{doṅgiliñcu}
  \item \emph{Say it:} \emph{doṅgiliñcu}, once more
  \item \emph{Say it:} \emph{doṅga nā sañci doṅgiliñcāḍu.} — The thief stole my bag.
  \item \emph{Recall:} say the Telugu for a war, then the Telugu for an incident, then say \emph{doṅgiliñcu} again
\end{itemize}

\begin{tcolorbox}[breakable,colback=teal!4,colframe=teal!35!black,title={Before you move on}]
What does \te{దొంగిలించు} mean? (\textquotedblleft{}To steal\textquotedblright{}.) Say it once more.
\end{tcolorbox}
\section[\texorpdfstring{\te{పారిపోవు} (pāripōvu)}{పారిపోవు (pāripōvu)}]{\te{పారిపోవు} (pāripōvu) — to run away; to escape}

\label{lesson:TE-C326-paripovu}



\begin{quote}
Before the new one: say the Telugu for an incident, then the Telugu for to steal.
\end{quote}

\subsection*{You'll want to know: \te{పారిపోవు}}
\textbf{\te{పారిపోవు}} — \emph{pāripōvu} — \textquotedblleft{}to run away; to escape\textquotedblright{}.

\begin{quote}
\te{పోలీసును} \te{చూసి} \te{దొంగ} \te{పారిపోయాడు}.
\end{quote}

\begin{quote}
\emph{pōlīsunu cūsi doṅga pāripōyāḍu.}
\end{quote}

\begin{quote}
Seeing the police, the thief ran away.
\end{quote}

Past (I …): \te{పారిపోయాను}, \emph{pāripōyānu}.

\te{పారి} is from \te{పారు}, to run; \te{పోవు} is to go. Many Telugu verbs end in -\te{పోవు} this way.

\begin{cousinweb}
\textbf{\te{పారిపోవు}} (\emph{pāripōvu}) is a native Dravidian compound: \emph{pāru} \textquotedblleft{}to run, run away\textquotedblright{} (DEDR 3963) and \emph{pōvu} \textquotedblleft{}to go\textquotedblright{} (DEDR 4572).
\end{cousinweb}

\begin{grammarlens}[title={What you've built}]
One more word for telling what happened.
\end{grammarlens}

\subsection*{Your turn}
\begin{itemize}
\raggedright
  \item \emph{Say it:} \emph{pāripōvu}
  \item \emph{Say it:} \emph{pāripōvu}, once more
  \item \emph{Say it:} \emph{pōlīsunu cūsi doṅga pāripōyāḍu.} — Seeing the police, the thief ran away.
  \item \emph{Recall:} say the Telugu for an incident, then the Telugu for to steal, then say \emph{pāripōvu} again
\end{itemize}

\begin{tcolorbox}[breakable,colback=teal!4,colframe=teal!35!black,title={Before you move on}]
What does \te{పారిపోవు} mean? (\textquotedblleft{}To run away; to escape\textquotedblright{}.) Say it once more.
\end{tcolorbox}
